

\medskip
\phantomsection\label{par:sga3-V-6-b}%
\noindent\textbf{b) Proof of \hyperref[lem:sga3-V-6.1-iii]{(iii)}.}---%
The groupoid \(V_\ast\) with base \(V\) is isomorphic both to the
inverse image of \(U_\ast\) under the base change \(u\) and to the
inverse image of \(X_\ast\) under the base change \(v\). We therefore
have a double Cartesian square:
\SGAThreeDiagram[0.65\textwidth]
  {figures/exp5/p272_lemma61_double_cartesian_square.png}
  {The double Cartesian square used in the proof of \hyperref[lem:sga3-V-6.1-iii]{Lemma~6.1(iii)}.}
  {fig:sga3-V-6.1-double-cartesian-square}
  {LNM 151, pp.~274--275}
Since \(u_0\boxtimes u_1\) is surjective, so is
\(v_0\boxtimes v_1\). Since \(v\times v\) is surjective, so are the
composite morphisms
\[
  V_1\xrightarrow{\ v_0\boxtimes v_1\ }V_0\times_YV_0
  \xrightarrow{\ v\times v\ }X_0\times_YX_0
  \qquad\text{and}\qquad
  V_1\longrightarrow X_1\xrightarrow{\ d_0\boxtimes d_1\ }X_0\times_YX_0,
\]
and hence so is \(d_0\boxtimes d_1\).

\medskip
\phantomsection\label{par:sga3-V-6-c}%
\noindent\textbf{c) Proof of \hyperref[lem:sga3-V-6.1-i]{(i)}.}---%
It remains to prove that \((Y,p)\) is a cokernel of \((d_0,d_1)\) in
the category of all ringed spaces. We first show that \(Y\) is
obtained from \(X_0\) by identifying those points \(x\) and \(y\) for
which there exists \(z\in X_1\) with
\[
  d_0(z)=x,\qquad d_1(z)=y.
\]
Indeed, \(p\) is surjective and \(p d_0=p d_1\). Conversely, if
\(p(x)=p(y)\), there is a point \(z'\in X_0\times_Y X_0\) whose first
projection is \(x\) and whose second projection is \(y\). If
\(z\in X_1\) satisfies
\[
  (d_0\boxtimes d_1)(z)=z',
\]
then \(d_0(z)=x\) and \(d_1(z)=y\), as required.

On the other hand, if \(W\) is a saturated open subset of \(X_0\),
then \(W\cap U\) is a saturated open subset of \(U\). By
Theorem~\hyperref[thm:sga3-V-4.1]{4.1}, \(q(W\cap U)\) is open in
\(Y\). Since \(q(W\cap U)=p(W)\), the topology on \(Y\) is the
quotient topology induced by \(X_0\).

It remains to show that the canonical sequence of sheaves of rings
\begin{equation}
  \mathcal O_Y
  \longrightarrow p_\ast\mathcal O_{X_0}
  \rightrightarrows
  p_\ast d_{0\ast}\mathcal O_{X_1}
  =
  p_\ast d_{1\ast}\mathcal O_{X_1}
  \label{eq:sga3-V-6.1-sheaf-exactness}
\end{equation}
is exact.
